\chapter{Algorithms and their cost}
\MARSChapterHeading{Algorithms and their cost}
\label{revision:algorithms-and-cost}

Two programs can produce the same ordered array while doing different amounts
of work. They may compare different pairs, move different words, and call the
display a different number of times. Choosing between them requires two separate
arguments: that each preserves the required result, and that the measurement
actually counts the work relevant to your question.

We will compare adjacent-exchange insertion sort with selection sort. Both grow
an ordered region of an array, but they choose the next operation differently.\footnote{For
  algorithm definitions and variants, see the NIST \emph{Dictionary of Algorithms and
    Data Structures} entries for insertion sort and selection sort:
  \url{https://xlinux.nist.gov/dads/HTML/insertionSort.html} and
  \url{https://xlinux.nist.gov/dads/HTML/selectionSort.html}.}
The complete public program is \texttt{sorting-costs.asm} in
\texttt{revision/labs/files}. It sorts up to eight signed words and draws an
initial frame, a frame after each complete exchange, and a final frame.

You should be able to state an ordering invariant, distinguish order from value
preservation, explain how equal keys behave, and compare measurements with an
explicit scope. The independent exercise changes the ordering requirement: group
negative words first while retaining the relative order inside both groups.

\section{Require both order and value preservation}

For an active prefix of length $n$, ascending order requires
$values[i]\leq values[i+1]$ whenever both positions exist. This condition alone is
not enough. Replacing every word with zero would satisfy it while losing the input.
The output must also contain the same values with the same multiplicities: it
must be a permutation of the active input. Words after that prefix are outside
the operation and must receive no stores.

The worked visualization admits $0\leq n\leq8$ and keys from $-4$ through three.
This small value domain gives each key a visible bar height from one through eight,
using $height=value+5$. It does not change signed order. Column $x$ is white at
screen row $y$ when $x<n$ and $y\geq3-values[x]$; other cells are black. Sorting
routines compare signed keys directly. The display's bounded height is a separate
reason for restricting this particular program's inputs.

Begin with the four-word prefix $[2,-2,3,-4]$. Both algorithms must produce
$[-4,-2,2,3]$. Table~\ref{tab:revision-sorting-exchanges} shows how their complete
exchanges differ. The full program's default has eight words; set its count to
four to follow this smaller example without touching the inactive suffix.

\begin{table}[htbp]
\centering
\begin{tabular}{lrl}
\toprule
Method & Exchanges so far & Active words\\
\midrule
Insertion & 0 & \texttt{2, -2, 3, -4}\\
 & 1 & \texttt{-2, 2, 3, -4}\\
 & 2 & \texttt{-2, 2, -4, 3}\\
 & 3 & \texttt{-2, -4, 2, 3}\\
 & 4 & \texttt{-4, -2, 2, 3}\\
\midrule
Selection & 0 & \texttt{2, -2, 3, -4}\\
 & 1 & \texttt{-4, -2, 3, 2}\\
 & 2 & \texttt{-4, -2, 2, 3}\\
\bottomrule
\end{tabular}
\caption{Complete exchanges on the same four-word input. A pass that makes no exchange produces no extra exchange snapshot. Both methods preserve the same multiset of values.}
\label{tab:revision-sorting-exchanges}
\end{table}


During an exchange, the first store can temporarily duplicate a value in memory
before the second restores the complete permutation. Publish the observation only
after both stores. A debugger can inspect either boundary, but a half-exchange is
not a completed sorting state.

\section{Extend a sorted prefix by adjacent exchanges}

Before insertion sort's outer iteration $i$, positions zero through $i-1$ are
already sorted. Start a cursor $j=i$ at the next item. Compare it with its left
neighbor. If it is strictly smaller, exchange the pair and move the cursor left;
otherwise its insertion position has been reached. The array up to $i$ is sorted
when this inner process ends, establishing the next outer iteration's invariant.

In the four-word trace, $-2$ first moves left across two. The next item, three,
already follows two and needs no exchange. Then $-4$ crosses three, two and $-2$,
one adjacent pair at a time. A display exchange is therefore not the same thing
as one completed outer iteration.

The following fragment belongs inside the complete \texttt{insertion\_sort}
function. Its setup saves the return link and \texttt{\$s0} through \texttt{\$s3},
sets \texttt{\$s0} to the array, \texttt{\$s1} to $n$, and \texttt{\$s2} to one.
The 48-byte frame contains five eight-byte saved slots and eight bytes of alignment
padding. The omitted epilogue restores those values and releases the frame.

\noindent\begin{minipage}{\linewidth}
  \MARSCodeLabel{Move the next item left until its prefix is ordered}
  \begin{minted}{mars-idea}
insertion_outer:
    bgec $s2, $s1, insertion_return
    nop
    move $s3, $s2
insertion_inner:
    beqzc $s3, insertion_next
    nop
    dsll $t0, $s3, 2
    daddu $t0, $s0, $t0
    lw $t1, 0($t0)
    lw $t2, -4($t0)
insertion_compare:
    bgec $t1, $t2, insertion_next
insertion_compare_end:
    nop
    sw $t1, -4($t0)
    sw $t2, 0($t0)
insertion_exchanged:
    move $a0, $s0
    move $a1, $s1
    balc publish_exchange
    addiu $s3, $s3, -1
    bc insertion_inner
insertion_next:
    addiu $s2, $s2, 1
    bc insertion_outer
\end{minted}
\end{minipage}

The zero-cursor check occurs before the left-neighbor load. When $j=0$, address
$base+4j-4$ would precede the active array, so there is no valid pair to compare.
The outer bound also handles empty and singleton prefixes without loading a key.
All address additions retain the full pointer width; the loaded words are sign
extended, and the compact comparison is signed.

An adjacent exchange of unequal out-of-order values removes exactly one inversion:
a pair of positions whose earlier key is greater than the later key. Thus this
implementation performs one exchange per initial inversion. An already sorted
prefix needs none. A reverse-ordered array of $n$ distinct keys has
$n(n-1)/2$ such pairs, giving 28 exchanges at $n=8$.

Equal keys do not exchange here. They therefore keep their relative order through
the adjacent moves: this insertion policy is stable. If records also had associated
payloads, those payloads would need to move with their keys. Our word-only array
does not store an independent identity for two equal integers.

\section{Choose the next minimum from the suffix}

Selection sort maintains a different invariant. Before iteration $i$, the first
$i$ positions contain the smallest $i$ keys in sorted order. Scan the remaining
suffix, retain the first occurrence of its minimum, and exchange that item with
position $i$ if the indices differ. This fixes one final position per outer iteration.

The complete function saves \texttt{\$ra} and \texttt{\$s0} through \texttt{\$s2}
in a 32-byte frame. They retain the array, $n$ and the left boundary across display
calls. Inside the scan, \texttt{\$t0} holds the minimum index, \texttt{\$t2} its
value, and \texttt{\$t1} the next index to inspect. These temporaries need not
survive a call because no observation call occurs until the chosen exchange is
complete.

\noindent\begin{minipage}{\linewidth}
  \MARSCodeLabel{Retain the first strictly smaller minimum}
  \begin{minted}{mars-idea}
selection_outer:
    addiu $t1, $s2, 1
    bgec $t1, $s1, selection_return
    nop
    move $t0, $s2
    dsll $t3, $s2, 2
    daddu $t3, $s0, $t3
    lw $t2, 0($t3)
selection_scan:
    bgec $t1, $s1, selection_exchange
    nop
    dsll $t3, $t1, 2
    daddu $t3, $s0, $t3
    lw $t4, 0($t3)
selection_compare:
    bgec $t4, $t2, selection_scan_next
selection_compare_end:
    nop
    move $t2, $t4
    move $t0, $t1
selection_scan_next:
    addiu $t1, $t1, 1
    bc selection_scan
\end{minted}
\end{minipage}

A candidate equal to the current minimum takes the branch and leaves the retained
index unchanged. This implements the stated first-minimum policy. When the scan
ends, the saved minimum value goes to the suffix's first position and that
position's old value goes to the minimum's old position. The complete source
skips the exchange if those indices already agree, then advances the boundary.

The scan compares $n-1$ candidates on the first iteration, $n-2$ on the next, and
so on. Its total is $n(n-1)/2$ key comparisons, regardless of the initial order.
It needs at most $n-1$ exchanges, and sometimes fewer because self-exchanges are
skipped.

Keeping the first minimum does not make a long-distance exchange stable. Give
equal keys conceptual labels: $[2_a,2_b,1]$. Selection's first exchange produces
$[1,2_b,2_a]$, reversing the two equal keys. The labels are a reasoning aid, not
extra bits stored in this program's array. Stability is a separate requirement
from ascending order and preservation of the key values.

\section{Preserve algorithm state across observations}

The shared \texttt{publish\_exchange} helper sets phase to one, increments the
exchange count, and calls \texttt{render}. It saves its own return link in a
16-byte frame. Both sorting functions call it only after completing their pair
of array stores. Their saved registers retain the loop state that must survive;
allowed temporary-register changes cannot corrupt the next algorithm step.

The renderer accepts an array pointer in \texttt{\$a0} and a valid prefix length
in \texttt{\$a1}. It reads only that prefix, writes all 64 framebuffer words,
preserves saved registers and the stack pointer, and reaches \texttt{sort\_ready}
before returning. It does not change the source, phase or exchange count. Its
caller owns publication metadata, just as the Life caller owned its generation
number. A separate renderer test can pass a different prefix address without
recreating the sorter's private register setup.

Main publishes phase zero with exchange count zero before sorting. After the
chosen routine returns, it publishes phase two with the final count. Empty,
singleton and already sorted arrays still have initial and final observations.
The frame must explicitly erase black and inactive cells on every render; otherwise
an earlier bar can remain visible after the array has changed.

Use \texttt{algorithm=0} for insertion and one for selection. At a
\texttt{sort\_ready} breakpoint, inspect phase, \texttt{swap\_count}, the active
words and the complete frame together. In the default eight-word input, insertion
has 17 exchanges and selection has five. A fast Run may repaint only the final
image, so visible screen refreshes are not an exchange counter.

\section{Measure a question with a stated scope}

A key comparison, a word store and a completed instruction are different units.
For the measurements below, the core records completed MARS instructions, including
expanded pseudoinstructions and executed \texttt{nop}s. One trace covers the entire
run. Another counts addresses from the chosen sorting routine's entry to its end
label. The latter includes its setup, comparisons, loop control, argument preparation,
calls and return, but excludes the bodies of \texttt{publish\_exchange} and
\texttt{render}. A third trace counts only the native branch instruction between
the routine's comparison labels, giving actual key comparisons for this source.

\begin{table}[htbp]
\centering
\small
\begin{tabular}{llrrrr}
\toprule
Input & Method & Key comparisons & Exchanges & Routine instr. & Whole instr.\\
\midrule
Sorted & Insertion & 7 & 0 & 101 & 2135\\
 & Selection & 28 & 0 & 317 & 2353\\
Reversed & Insertion & 28 & 28 & 479 & 30933\\
 & Selection & 28 & 4 & 409 & 6505\\
Mixed & Insertion & 22 & 17 & 344 & 19633\\
 & Selection & 28 & 5 & 399 & 7510\\
\bottomrule
\end{tabular}
\caption{Observed completed-instruction counts for the exact worked source on MARS64 Release 6. Routine intervals exclude the publication and renderer bodies; whole-run counts include them. Each input contains all eight keys from minus four through three.}
\label{tab:revision-sorting-costs}
\end{table}


The three inputs use the same eight distinct keys. Their order changes, while
the displayed heights and total number of white cells stay the same. Each method
starts from a fresh copy of its input. Whole-run counts include the common caller,
metadata, initial and final frames, exchange frames and exit. Assembly, session
construction and host time are outside all these counts.

The mixed row demonstrates why scope matters. Insertion uses 344 instructions
inside its sorting routine, compared with selection's 399. Across the whole
visualized program, insertion uses 19,633 and selection 7,510. Insertion's 17
exchanges cause 19 complete renders, including initial and final; selection's five
cause seven. That is 1,216 versus 448 framebuffer word stores. The observation
work changes the result of the comparison.

These are measurements of these implementations, not hardware cycles or universal
algorithm timings. The routine interval still contains observation-call preparation
and calls, so it is not a claim to count only comparisons and moves. Removing the
visualizer, changing its implementation, or changing the calling structure creates
a different measured program. Measure the workload you intend to run.

Repeat the measurements with JDK 25:
\begin{verbatim}
python3 tools/measure_revision_sorting.py \
  --java /path/to/jdk-25/bin/java \
  --cli /path/to/cli.jar --output /new/report/directory
\end{verbatim}
The tool checks actual arrays, metadata, frame contents and store footprints while
recording the core traces. Its report binds the source and CLI hashes. Compare
those bindings before treating a changed number as a regression or improvement.

The formulas also explain growth beyond this small demonstration. Selection scans
a quadratic number of key pairs. Insertion has a quadratic worst case, but an
already sorted input needs only one key comparison per new item. Those statements
explain input sensitivity and scaling; they do not replace the concrete measurements
or establish host execution time.

\section{Apply an ordering requirement independently}

The separate \texttt{partition-negative.asm} exercise groups negative words before
nonnegative words while preserving relative order inside each group. For example,
$[2,-2,3,-4,0,1,-3,-1]$ must become
$[-2,-4,-3,-1,2,3,0,1]$, with a returned negative count of four. The negative group
is not required to be numerically sorted.

This task accepts the full signed 32-bit word range, unlike the visualizer's small
bar-height domain. Zero belongs to the nonnegative group. Choose your own valid
algorithm: you might move a newly encountered negative word left through earlier
nonnegative words, or use private temporary storage and copy the two groups back.
Preserve inactive words, guards, saved registers and caller-owned stack state,
and limit source reads and array stores to the supplied active span.

State the invariant that preserves each group's order, predict empty and all-one-sign
cases, and test the most negative word as well as zero. The written lab supplies
public checks and counterexamples. Sorting the whole prefix into ascending order
would satisfy the sign boundary but violate the required relative order; use the
actual contract to decide whether a result is correct.
